\cvsection{Expériences professionnelles}

\begin{cventries}
  \cventry
    {Stagiaire ingénieur de recherche en IA}
    {INRIA}
    {Bordeaux, France}
    {Mars 2026 -- Septembre 2026}
    {
      \begin{cvitems}
        \item {Conception et développement en autonomie, avec le suivi de mes encadrants, d’une application fonctionnelle pour aider les chercheurs à trouver des articles scientifiques.}
        \item {Développement d’un backend Python/FastAPI avec API REST documentées avec OpenAPI et persistance des documents et métadonnées dans MongoDB.}
        \item {Développement d’une interface React/Next.js pour interagir avec l’agent IA et ses outils, avec sélection des messages à inclure dans le contexte du LLM.}
        \item {Intégration de l’API OpenAI dans un agent multi-étapes et transmission des résultats de tâches cron par webhooks.}
        \item {Versionnement avec Git/GitHub, conteneurisation des composants avec Docker et mise en place d’une CI/CD.}
        \item {Intégration d’un pipeline RAG combinant MongoDB et ChromaDB, avec chunking, embeddings et recherche sémantique.}
        \item {Gestion d’une application utilisant des modèles exécutés localement avec LM Studio sur un Mac Studio, avec installation et configuration des modèles.}
        \item {Évaluation de la qualité et de la latence des LLM sur HPC avec GPU NVIDIA et vLLM, analyse des résultats et rédaction en cours d’un article scientifique.}
      \end{cvitems}
    }

  \cventry
    {Stagiaire ingénieur de recherche en apprentissage par renforcement}
    {Université Karunya}
    {Coimbatore, Inde}
    {Juin 2025 -- Sept. 2025}
    {
      \begin{cvitems}
        \item {Développement d’un modèle d’apprentissage par renforcement et de son environnement d’entraînement pour la navigation autonome d’un robot hospitalier.}
        \item {Déploiement et optimisation du logiciel et du modèle sur Raspberry Pi pour des tests en conditions réelles.}
      \end{cvitems}
    }
\end{cventries}
